\documentclass{article}
\usepackage{amsmath}
\usepackage{gensymb}
\usepackage{amsfonts}
\usepackage{amssymb}
\usepackage{geometry} % Required for customizing margins

\geometry{margin=1in} % Sets 1-inch margins on all sides


\title{EMCH 360 Assignment \#5 \& 6}
\author{Jacob C. Vaught}
\date{December 1st 2023}

\begin{document}

\maketitle
\section*{Problem \#1}
Given a 1:15 scale model of a submarine to be tested at 180 ft/s in a wind tunnel with standard sea-level air, determine the speed of the prototype in seawater ensuring Reynolds number similarity. The kinematic viscosities at 60\degree F are given as:

\begin{itemize}
    \item Water: $1.21 \times 10^{-5}$ ft\(^2\)/s
    \item Air: $1.58 \times 10^{-4}$ ft\(^2\)/s
\end{itemize}

Applying Reynolds number similarity:

\subsection*{Solution}
The Reynolds number similarity is given by the equation:
\[ \frac{V_{\text{model}} D_{\text{model}}}{\nu_{\text{air}}} = \frac{V_{\text{prototype}} D_{\text{prototype}}}{\nu_{\text{water}}} \]

Given the scale of the model, \( D_{\text{prototype}} = 15 D_{\text{model}} \). Rearranging the equation to find \( V_{\text{prototype}} \):

\[ V_{\text{prototype}} = V_{\text{model}} \times \frac{\nu_{\text{water}}}{\nu_{\text{air}}} \times \frac{D_{\text{model}}}{D_{\text{prototype}}} \]

Substituting the given values:

\begin{align*}
    V_{\text{prototype}} &= 180 \, \text{ft/s} \times \frac{1.21 \times 10^{-5} \, \text{ft}^2/\text{s}}{1.58 \times 10^{-4} \, \text{ft}^2/\text{s}} \times \frac{1}{15} \\
    &= 180 \, \text{ft/s} \times \frac{1.21 \times 10^{-5}}{1.58 \times 10^{-4}} \times \frac{1}{15} \\
    &\approx 0.919 \, \text{ft/s}
\end{align*}

Hence, the speed of the prototype to ensure Reynolds number similarity is approximately 0.919 ft/s.

\newpage


\section*{Problem \#2}

Given the pump power $P = 200 \frac{ft \cdot lb}{s}$, pipe length $L = 200$ ft, volumetric flow rate $Q = 0.05 \frac{ft^3}{s}$, specific weight of water $\omega = 62.4 \frac{lb}{ft^3}$, and gravitational acceleration $g = 32.2 \frac{ft}{s^2}$, the solution is as follows:

\subsection*{Part (a): Pump Head Development}
The head developed by the pump $h_p$ is calculated using the power equation $P = \omega Q h_p$, resulting in:
\[ h_p = \frac{P}{\omega Q} = \frac{200}{62.4 \times 0.05} = 64.102 \text{ ft}. \]

\subsection*{Part (b): Major Loss Calculation}
For the pipe diameter $d = 0.1$ ft, the velocity $v$ and Reynolds number $Re$ are:
\[ v = \frac{Q}{\frac{\pi d^2}{4}} = \frac{0.05}{\frac{\pi(0.1)^2}{4}} = 6.366 \frac{ft}{s}, \]
\[ Re = \frac{vd}{\nu} = \frac{6.366 \times 0.1}{1.08 \times 10^{-5}} = 58,944.44. \]

Using the friction factor $f = 0.039$ from the Moody chart for $Re = 58,944.44$ and $\frac{\epsilon}{d} = 0.01$, the major loss is:
\[ h_{\text{major}} = f \frac{L}{D} \frac{v^2}{2g} = 0.039 \frac{200}{0.1} \frac{6.366^2}{2 \times 32.2} = 49.084 \text{ ft}. \]

\subsection*{Part (c): Minor Loss Calculation}
The minor losses $h_{\text{minor}}$ are calculated as:
\[ h_{\text{minor}} = (\Sigma K_L) \frac{v^2}{2g} = (0.8 + 1.2 + 6 + 1 + 1.5) \times \frac{6.366^2}{2 \times 32.2} = 13.403 \text{ ft}. \]

\subsection*{Summary}
The final answers are:
\begin{itemize}
    \item Head developed by the pump: $64.102$ ft.
    \item Major loss: $49.084$ ft.
    \item Minor loss: $13.403$ ft.
\end{itemize}

\newpage
\section*{Problem \#3}
From the continuity equation, the velocity \( V_2 \) is given by:
\[ V_2 = \frac{Q}{A} = \frac{Q}{\frac{\pi d^2}{4}} = \frac{0.1}{\frac{\pi (0.125)^2}{4}} = 8.14 \frac{ft}{s} \]

Applying the Bernoulli equation between sections 1 and 2 and assuming \( P_1 = 0 \), we get:
\[ P_1 = \left( f \frac{L}{d} + \Sigma K_L \right) \frac{1}{2} \rho V^2 \]

Where \( \rho = 1.53 \times 10^{-3} \frac{slugs}{ft^3} \) and \( \mu = 4.73 \times 10^{-7} \frac{lb \cdot s}{ft^2} \).

The relation between the surface roughness and diameter is \( \frac{\epsilon}{d} = 0.0068 \) for cast iron.

The Reynolds number \( Re \) is calculated as:
\[ Re = \frac{\rho V d}{\mu} = \frac{(1.53 \times 10^{-3}) \times 8.14 \times 0.125}{4.73 \times 10^{-7}} = 3,291 \]

From the Moody chart, for \( Re = 3,291 \) and \( \frac{\epsilon}{d} = 0.0068 \), the friction factor \( f \) is 0.047.

Substituting all values into the pressure equation:
\[ P_1 = \left( 0.047 \frac{14}{0.125} + (6 \times 0.3 + 8.5) \right) \frac{1}{2} \times (1.53 \times 10^{-3}) \times (8.14)^2 \]
\[ P_1 = 0.9 \frac{lb}{ft^2} \]





\newpage
\section*{Problem \#4}

Let's apply the Bernoulli equation considering the losses from entry (1) to exit (2), assuming $v_1 = v_2$ and $z_1 = z_2$:

\[ p_1 - p_2 = \rho g h_l \]

where $h_l$ is the total headloss in pipe flow given by:

\[ h_l = \frac{v^2}{2g} \left( f \frac{L}{d} + nk \right) \]

which implies the pressure drop $\Delta p$ is:

\[ \Delta p = \rho \frac{v^2}{2} \left( f \frac{L}{d} + nk \right) \]

From the table of the roughness values of different pipes, for copper pipe (drawn tubing), we get $\epsilon = 5 \times 10^{-6} \text{ ft}$.

From the properties table of water at $40^\circ F$ we have $\nu = 1.66 \times 10^{-5} \text{ ft}^2/\text{s}$ and $\rho = 1.94 \text{ slugs/ft}^3$.

Let's find the velocity first: $v = \frac{\text{discharge}}{\text{Area}}$

\[ V = \frac{0.9 \text{ gal/min}}{\frac{\pi}{4} d^2} = \frac{0.9 \times 0.133681 \text{ ft}^3/\text{min}}{60 \times \frac{\pi}{4} \left( \frac{0.5}{12} \right)^2 \text{ ft}^2} \approx 1.47059 \text{ ft/s} \]

Now, let's find Reynolds's number:

\[ Re = \frac{\rho V d}{\nu} = \frac{1.94 \times 1.47059 \times \left( \frac{0.5}{12} \right)}{1.66 \times 10^{-5}} \approx 3,691.24 \]

The relative roughness $\frac{\epsilon}{d}$ is:

\[ \frac{\epsilon}{d} = \frac{5 \times 10^{-6}}{\left( \frac{0.5}{12} \right)} = 1.2 \times 10^{-4} \]

Using the friction factor from the Moody chart:

\[ f = 0.041 \]

We have:

\[ \Delta p = \rho \frac{v^2}{2} \left( f \frac{L}{d} + nk \right) \]

Putting the values we get:

\[ \Delta p = 1.94 \frac{(1.47059)^2}{2} \left( 0.041 \times \frac{12}{\left( \frac{0.5}{12} \right)} + 7 \times 1.5 \right) = 45.4207 \text{ lb/ft}^2 \]

Converting to psi:

\[ \Delta p = \frac{45.4207}{144} = 0.315 \text{ psi} \]

\textbf{Final answer:}

\[ \Delta p = 0.315 \text{ psi} \]

\newpage
\section*{Poblem \#5}

Apply Bernoulli's equation between points 1 and 2 located at water surface elevation of lower and upper tanks respectively,
\[ \frac{P_1}{\gamma} + \frac{V_1^2}{2g} + Z_1 + H_p = \frac{P_2}{\gamma} + \frac{V_2^2}{2g} + Z_2 + H_L, \]
where $P$ is pressure, $V$ is velocity, $Z$ is elevation, $H_p$ is head added by the pump, and $H_L$ is head loss in the system.

Since both points are open to the atmosphere, pressure is only atmospheric so gauge pressure $P_1 = P_2 = 0$ and velocity at the surface is negligible so $V_1 = V_2 = 0$. Elevation, $Z_1 = 10$ m and $Z_2 = 13$ m.

The above equation is reduced to,
\[ 0 + 0 + 10 + H_p = 0 + 0 + 13 + H_L, \]
\[ \Rightarrow H_p = 3 + H_L \dots (1) \]

Now, Major head loss due to friction ($H_L$) is given by,
\[ H_L = f \frac{L V^2}{2gD}, \]
where $L$ is pipe length, $f$ is friction factor, $V$ is flow velocity, $D$ is pipe diameter, and $g$ is acceleration due to gravity.

Given:
\begin{itemize}
    \item Pipe length $L = 80$ m.
    \item Flow velocity $V = \frac{Q}{A} = \frac{0.1}{\frac{\pi}{4}(0.15)^2} = 5.66$ m/s.
    \item Pipe diameter $D = 15$ cm or $0.15$ m.
    \item Kinematic viscosity of water $\nu = 10^{-6}$ m$^2$/s.
    \item Roughness of forged steel $e = 0.045$ mm.
\end{itemize}

Reynolds number ($Re$) is calculated as,
\[ Re = \frac{VD}{\nu} = \frac{5.66 \times 0.15}{10^{-6}} = 6.5 \times 10^5. \]

Relative roughness ($\frac{e}{D}$) is,
\[ \frac{e}{D} = 0.0145. \]

From the Moody diagram, for $Re = 8.49 \times 10^5$ and $\frac{e}{D} = 0.0003$, the friction factor ($f$) is approximately $0.0156$. Thus, the head loss due to friction is,
\[ H_L = 0.0145 \frac{80 \times 5.66^2}{2 \times 9.81 \times 0.15} = 12.614 \text{ m}. \]

Putting values in equation (1),
\[ H_p = 3 + 12.614 = 15.614 \text{ m}. \]

Therefore, the power added by the pump ($P$) is,
\[ P = \gamma Q H_p, \]
where $\gamma$ is the unit weight of water $= 9810$ N/m$^3$.

\[ P = 9810 \times 0.1 \times 15.614, \]
\[ P = 16270 \text{ W} \text{ or } 16.27 \text{ kW}. \]

\textbf{Final answer:}
The power added to water by the pump is $15.61$ kW.

\newpage
\section*{Problem \#6}

\subsection*{Given:}
\begin{itemize}
    \item The velocity of the water at the dam model \( v_m = 5 \) m/s.
    \item The scale of the model is \( s = 1:15 \).
\end{itemize}

\subsection*{Explanation:}
This question is solved by use of dimensional analysis. For the given conditions, Froude's number is used, i.e., the Froude's number must be the same for both the model and the dam.

\subsection*{Calculation:}
The expression for the Froude's number is
\[ Fr = \frac{v}{\sqrt{g \cdot l}} \]
where \( v_m \), \( g_m \), \( l_m \) is the velocity of water at the model, acceleration due to gravity at the model, and characteristic dimension of the model respectively, and \( v \), \( g \), \( l \) is the velocity of water at the dam, acceleration due to gravity at the dam, and characteristic dimension of the dam respectively.

Since the Froude's number must be consistent across both scales and gravity is constant, we have
\[ Fr_m = Fr \Rightarrow \frac{v_m}{\sqrt{g \cdot l_m}} = \frac{v}{\sqrt{g \cdot l}} \]

Substituting the value \( s = \frac{1}{15} \) for the scale ratio and \( v_m = 5 \) m/s for the velocity of the model, we get
\[ v = v_m \cdot \sqrt{\frac{1}{s}} \]
\[ v = 5 \cdot \sqrt{15} \]
\[ v \approx 19.36 \text{ m/s} \]

\subsection*{Final Answer:}
The velocity of water at the prototype is approximately 19.36 m/s.

\newpage
\section*{Problem \#7}

Given:
\begin{itemize}
    \item Viscosity of the oil, \( \mu = 0.4 \) N.s/m\(^2\).
    \item Diameter of the pipe, \( D = 0.03 \) m.
    \item Length of the pipe, \( L = 15 \) m.
    \item Flow rate, \( Q = 3.0 \times 10^{-3} \) m\(^3\)/s.
\end{itemize}

The pressure drop \( \Delta P \) across a horizontal pipe can be calculated using Poiseuille's law, which for a circular pipe is:

\[ \Delta P = \frac{128 \mu L Q}{\pi D^4} \]

Substituting the given values into the equation:

\[ \Delta P = \frac{128 \times 0.4 \times 15 \times 3.0 \times 10^{-3}}{\pi \times (0.03)^4} \]
\[ \Delta P \approx 603609.858 \text{ Pa} \]

\textbf{Final Answer:}

The pressure drop across the pipe is approximately \( 338.6 \) MPa.
\newpage



\section*{Problem \#8}
A 2-meter diameter pipe made of finished concrete (Manning roughness coefficient, \( n = 0.012 \)) lies on a slope of 1-meter elevation change per 1000-meter horizontal distance. The goal is to determine the flow rate when the pipe is half full.

\section*{Calculation}
The Manning formula for open channel flow is given by:
\begin{equation}
    Q = \frac{1}{n} A R^{\frac{2}{3}} S^{\frac{1}{2}}
\end{equation}
where:
\begin{itemize}
    \item \( Q \) is the flow rate in cubic meters per second (m\(^3\)/s).
    \item \( n \) is the Manning roughness coefficient.
    \item \( A \) is the cross-sectional area of flow in square meters (m\(^2\)).
    \item \( R \) is the hydraulic radius in meters (m), equal to the area of flow divided by the wetted perimeter.
    \item \( S \) is the slope of the pipe.
\end{itemize}

Given that the diameter of the pipe, \( D \), is 2 meters, the radius, \( r \), is 1 meter. The cross-sectional area \( A \) when the pipe is half full is:
\begin{equation}
    A = \frac{1}{2} \pi \times (1\, \text{m})^2 = 1.5708\, \text{m}^2
\end{equation}

The wetted perimeter \( P \) is:
\begin{equation}
    P = \pi \times 1\, \text{m} + 2\, \text{m} = 5.1416\, \text{m}
\end{equation}

The hydraulic radius \( R \) is:
\begin{equation}
    R = \frac{1.5708\, \text{m}^2}{5.1416\, \text{m}} = 0.3058\, \text{m}
\end{equation}

Substituting these values into the Manning formula, we get the flow rate \( Q \):
\begin{equation}
    Q = \frac{1}{0.012} \times 1.5708\, \text{m}^2 \times (0.3058\, \text{m})^{\frac{2}{3}} \times (1/1000)^{\frac{1}{2}} \approx 1.878\, \text{m}^3/\text{s}
\end{equation}

\section*{Results}
After performing the calculations, the flow rate \( Q \) is found to be approximately 1.878 m\(^3\)/s.
\newpage
\section*{Problem \#9}
Given an earthen channel with a bottom width of 4 ft and a top width of 12 ft at a depth of 3 ft, the task is to find the minimum slope required to maintain a flow velocity of at least 2.5 ft/s, using the Manning formula with a roughness coefficient \( n \) of 0.022 and a conversion factor \( k \) of 1.49.

\section*{Calculation}
The Manning's formula in US customary units is given by:
\begin{equation}
    v = \frac{k}{n} R^{\frac{2}{3}} S^{\frac{1}{2}}
\end{equation}
where:
\begin{itemize}
    \item \( v = 2.5 \) ft/s is the desired flow velocity.
    \item \( k = 1.49 \) is the conversion factor for US customary units.
    \item \( n = 0.022 \) is the Manning roughness coefficient.
    \item \( R \) is the hydraulic radius in feet.
    \item \( S \) is the slope of the channel bottom.
\end{itemize}

The hydraulic radius \( R \) is calculated as the area of flow (\( A \)) divided by the wetted perimeter (\( P \)). For the given trapezoidal channel, the area \( A \) and the wetted perimeter \( P \) are calculated as follows:
\begin{align}
    A &= 4\, \text{ft} \times 3\, \text{ft} + 3\, \text{ft} \times \left(\frac{12\, \text{ft} - 4\, \text{ft}}{2}\right) = 30\, \text{ft}^2 \\
    P &= 4\, \text{ft} + 2 \times \sqrt{3\, \text{ft}^2 + \left(\frac{12\, \text{ft} - 4\, \text{ft}}{2}\right)^2} = 17.464\, \text{ft}
\end{align}

The hydraulic radius \( R \) is therefore:
\begin{equation}
    R = \frac{30\, \text{ft}^2}{17.464\, \text{ft}} = 1.718\, \text{ft}
\end{equation}

The slope \( S \) can be rearranged from Manning's formula to find the minimum value required to achieve the desired velocity:
\begin{equation}
    S = \left(\frac{2.5\, \text{ft/s} \cdot 0.022}{1.49 \cdot 1.718\, \text{ft}^{\frac{2}{3}}}\right)^2 = 0.000664
\end{equation}

Therefore, the minimum slope \( S \) is calculated to be approximately 0.000664, or 0.0664\%.



\newpage
\section*{Problem \#10}

Given the following parameters for water flow over a sharp-crested weir:
\begin{itemize}
    \item Width of the weir, \( b = 2 \) m.
    \item Head over the weir, \( h = 0.1 \) m.
    \item Upstream channel depth, \( d = 1 \) m (not required in the formula, but given for completeness).
    \item Acceleration due to gravity, \( g = 9.81 \) m/s\(^2\).
    \item Discharge coefficient for a sharp-crested weir, \( C_d \approx 0.611 \).
\end{itemize}

The flow rate \( Q \) can be determined using the following equation:
\[ Q = \frac{2}{3} C_d b h^{3/2} \sqrt{2g} \]

Substituting the given values into the equation, we have:
\[ Q = \frac{2}{3} \times 0.611 \times 2 \times 0.1^{3/2} \times \sqrt{2 \times 9.81} \]

Calculating the flow rate, we get:
\[ Q = \frac{2}{3} \times 0.611 \times 2 \times 0.1^{1.5} \times \sqrt{19.62} \]
\[ Q = \frac{2}{3} \times 0.611 \times 2 \times 0.03162 \times 4.429 \]
\[ Q = \frac{2}{3} \times 0.611 \times 2 \times 0.14005 \]
\[ Q = \frac{2}{3} \times 0.611 \times 0.2801 \]
\[ Q = \frac{2}{3} \times 0.17115 \]
\[ Q = 0.1141 \text{ m}^3/\text{s} \]

Thus, the flow rate over the weir is approximately \( 0.1141 \) cubic meters per second.


\end{document}
